\subsection{Rechenregeln}
\label{subsec:vector_algebra_scalar_product_rules}

\begin{TheoremBox}{Rechenregeln des Skalarprodukts}{vector_algebra_dot_product_laws}
  Für $\vec a,\vec b,\vec c\in\R^n$ und $\lambda\in\R$ gilt
  \[
    \begin{aligned}
      \vec a\cdot\vec b
        &=\vec b\cdot\vec a
        &&\text{(Kommutativgesetz)},\\[0.5em]
      \vec a\cdot(\vec b+\vec c)
        &=\vec a\cdot\vec b+\vec a\cdot\vec c
        &&\text{(Distributivgesetz)},\\[0.5em]
      (\lambda\vec a)\cdot\vec b
        &=\lambda(\vec a\cdot\vec b)
        =\vec a\cdot(\lambda\vec b).
    \end{aligned}
  \]
\end{TheoremBox}

\begin{RemarkBox}[title=Keine Assoziativität]
  \label{rem:vector_algebra_dot_product_not_associative}
  Ein Assoziativgesetz ist nicht möglich: $\vec a\cdot\vec b$ ist bereits
  ein Skalar. Daher sind $(\vec a\cdot\vec b)\cdot\vec c$ und
  $\vec a\cdot(\vec b\cdot\vec c)$ keine Skalarprodukte zweier Vektoren.
\end{RemarkBox}

\begin{ExampleBox}{Rechenregeln anwenden}{vector_algebra_dot_product_laws_example}
  Für
  \[
    \vec a=\begin{pmatrix}1\\2\end{pmatrix},\qquad
    \vec b=\begin{pmatrix}3\\-1\end{pmatrix},\qquad
    \vec c=\begin{pmatrix}-1\\4\end{pmatrix}
  \]
  folgt
  \[
    \vec a\cdot(\vec b+\vec c)
    =\vec a\cdot\vec b+\vec a\cdot\vec c
    =\Cloze{5cm}{1+7=8}.
  \]
\end{ExampleBox}
